\chapter{Пространственно-временное отношение кинетики и Hodge-потенциала}
\label{ch:v7-spacetime-kinetic-potential-ratio-admission-gate}

\section{Произведённый оператор}

Одна массовая калибровка оставила свободной эффективную квартику
$\lambda_E=\kappa/Z^2$. Проверим, сокращается ли эта свобода в минимальном
плоском произведённом операторе
\begin{equation}
 \mathcal D_E
 =i\gamma^\mu\partial_\mu\otimes I
 +\gamma^5\otimes a\Phi_E(x),
 \label{eq:v7-spacetime-ratio-product-dirac}
\end{equation}
где $a>0$ является общей нормировкой нечётного поля, а $\Phi_E$ содержит
динамический и фоновый цвета Hodge-суперсвязности.

В плоском четырёхмерном фоне $a_4$ даёт кинетический и квартичный члены с
одним коэффициентом \cite{v7-spacetime-ratio-spectral-action}:
\begin{equation}
 C_0=\frac{f_0}{8\pi^2}.
 \label{eq:v7-spacetime-ratio-common-coefficient}
\end{equation}
Квадратичный уровень вакуума зависит также от $a_2$, но после одной
массовой калибровки он не входит в искомую безразмерную квартику.

\section{Конечные следовые множители}

Для динамического эрмитова оператора рёбер
$\Phi_Z=d_Z+d_Z^\dagger$ выполняется
\begin{equation}
 \Tr_{\mathcal K_E}(\partial_\mu\Phi_Z)^2
 =2\sum_e\left[(\partial_\mu x_e)^2+(\partial_\mu y_e)^2\right].
 \label{eq:v7-spacetime-ratio-kinetic-trace}
\end{equation}
Hodge-потенциал был определён равенством
\begin{equation}
 \Tr_{\mathcal K_E}\mathfrak m_\mu^2
 =2\mathcal S_\mu+10\mu^4.
 \label{eq:v7-spacetime-ratio-potential-trace}
\end{equation}
Поэтому произведённый $a_4$-блок индуцирует коэффициенты
\begin{equation}
 Z=4C_0a^2,
 \qquad
 \kappa=2C_0a^4.
 \label{eq:v7-spacetime-ratio-induced-coefficients}
\end{equation}

\section{Что действительно сокращается}

После канонической нормировки
\begin{equation}
 \boxed{
 \lambda_E=\frac{\kappa}{Z^2}
 =\frac1{8C_0}=\frac{\pi^2}{f_0}.}
 \label{eq:v7-spacetime-ratio-effective-quartic}
\end{equation}
Общая нормировка $a$ сокращается точно. Аналогично единый массовый масштаб
равен
\begin{equation}
 M_0^2=\frac{\kappa\mu^2}{Z}
 =\frac{a^2\mu^2}{2},
 \label{eq:v7-spacetime-ratio-mass-scale}
\end{equation}
поэтому $f_0$ сокращается из линейных масс и одна калибровка фиксирует
комбинацию $a\mu$.

\begin{theorem}[Сокращение нормировки нечётного поля]
В минимальном произведённом heat-kernel блоке отношение
$\kappa/Z^2$ не зависит от рескейлинга $\Phi_E\mapsto a\Phi_E$, от
$f_2$ и от cutoff после фиксации одной массы. Вся оставшаяся
неопределённость эффективной квартики сосредоточена в одном безразмерном
моменте $f_0$.
\end{theorem}

\section{Почему численное значение ещё не предсказано}

Функция спектрального отсечения определяет $f_0$ и не выводится текущим
конечным графом. При
\begin{equation}
 f_0\in\left\{\frac12,1,2,5\right\}
 \quad\Longrightarrow\quad
 \lambda_E\in\left\{2\pi^2,\pi^2,\frac{\pi^2}{2},\frac{\pi^2}{5}\right\}.
 \label{eq:v7-spacetime-ratio-f0-family}
\end{equation}
Все эти значения сохраняют прежнее отношение линейных масс $\sqrt2$, но
дают разные амплитуды рассеяния.

Полное Real-удвоение умножает оба следовых члена на два. Физический
полуслед удаляет этот известный множитель; без него каноническая квартика
изменилась бы вдвое. Семейная и Real-кратности поэтому должны фиксироваться
тем же физическим следом, а не выбираться после нормировки.

Обычный выход спектрального действия состоит в фиксации $f_0$ через
калибровочный кинетический член. Но текущий вспомогательный рёберный носитель
не содержит доказанного общего физического gauge-блока. До такого вложения
формула~\eqref{eq:v7-spacetime-ratio-effective-quartic} является отношением
к одному невыведенному моменту, а не численным предсказанием. Произведённые
спектральные тройки дают подходящую архитектуру для следующей проверки
\cite{v7-spacetime-ratio-product}.

\begin{nogocriterion}[Один оставшийся безразмерный якорь]
Нельзя назначать $f_0=1$ или определять его по желаемой квартике. Полное
одномасштабное замыкание допускается только если один и тот же физический
спектральный след фиксирует калибровочную кинетику и
$\lambda_E=\pi^2/f_0$. Без общего gauge-блока численное значение
$\lambda_E$ остаётся свободным.
\end{nogocriterion}

\begin{protocol}
\begin{enumerate}
 \item общий коэффициент $a_4$: \textbf{$C_0=f_0/(8\pi^2)$};
 \item кинетический коэффициент: \textbf{$Z=4C_0a^2$};
 \item потенциальный коэффициент: \textbf{$\kappa=2C_0a^4$};
 \item эффективная квартика: \textbf{$\lambda_E=\pi^2/f_0$};
 \item зависимость от нормировки $a$: \textbf{сокращается};
 \item зависимость от $f_2$ и cutoff: \textbf{сокращается после
 массовой калибровки};
 \item зависимость от $f_0$: \textbf{остаётся};
 \item Real-полуслед: \textbf{обязателен и сохраняет коэффициенты};
 \item общий физический gauge-блок: \textbf{не выведен};
 \item итог: \textbf{частичный проход, свобода сведена к одному $f_0$}.
\end{enumerate}
\end{protocol}

Машинный сертификат сохранён в
\path{s2t_v7_spacetime_kinetic_potential_ratio_admission_gate_results.json}.

\begin{thebibliography}{9}
\bibitem{v7-spacetime-ratio-spectral-action}
A.~H.~Chamseddine, A.~Connes,
\emph{The Spectral Action Principle},
Communications in Mathematical Physics 186 (1997), 731--750;
arXiv:hep-th/9606001.
\bibitem{v7-spacetime-ratio-product}
S.~Brain, B.~Mesland, W.~D.~van Suijlekom,
\emph{Gauge Theory for Spectral Triples and the Unbounded Kasparov Product},
Journal of Noncommutative Geometry 10 (2016), 135--206;
arXiv:1306.1951.
\end{thebibliography}